\section{Existence and unique selection of the solved rate}
\label{sec:existence}

This section proves the two facts from \cref{sec:setup:roadmap}: $E$ is
non-increasing (Monotonicity), and its sign changes across a search range
whose endpoints we can write down (Bracketing). Together they give a unique
largest self-consistent rate, which bisection finds deterministically.

\begin{theorem}[the solved rate exists and is uniquely selected]
\label{thm:main}
Fix $\STraw \geq 0$, the checkpoint, the supply $N \geq 1$, the junior
share $f \leq 1$, and the frozen invariant $L$ at the start of a sync. On the bracket
\[
  \mathcal{B} \;=\; \bigl[\,0,\ x_{\max}\,\bigr],
  \qquad
  x_{\max} \;=\; \STraw + \lfloor f\,V_{\max}\rfloor + 1,
\]
where $V_{\max} = \sup_y V(y)$ is the pool value at the all-quote corner
(the senior shares all sold), the integer error function
$E(x) = F(P_A(x/N)) - x$ is
non-increasing with steps in $\{-1, 0\}$ \textup{(Monotonicity)} and
satisfies $E(0) \geq 0$ and
$E(x_{\max}) < 0$ \textup{(Bracketing)}. Hence the integers where $E \geq 0$
form an initial segment $[0, x^\star]$ with a unique right endpoint $x^\star$, the
largest self-consistent senior NAV, and bisection on the predicate
$E(x) \geq 0$ finds it.
\end{theorem}

\Cref{sec:existence:mono,sec:existence:sign} prove the two pillars, and
\cref{sec:existence:generality} records what the argument assumes about
the pool, so an auditor can check it per market.

\subsection{Monotonicity: $E$ is non-increasing}
\label{sec:existence:mono}

$E(x) = F(P_A(x/N)) - x$. The $-x$ term falls by one per step, so it is
enough to show the composition $x \mapsto F(P_A(x/N))$ rises by at most
one per step. We bound each stage in turn.

\paragraph{Stage 1: $f\widehat{V}$ moves by at most $1$ per step, so $P_A$ moves by $0$ or $1$.}
Fix the invariant $L$ at the pre-solve cached rate, as
\cref{sec:setup:PA} requires. The fair pool value at rate $y$ on the
frozen curve is the quote it holds plus its senior shares marked at $y$,
\[
  V(y) \;=\; b_q(y) \;+\; y\,D(y),
\]
where $D(y)$ is the pool's demand for senior shares at price $y$ and
$b_q(y)$ is the quote reserve at the same fair point. Moving along the
curve $db_q/dD = -y$ (the slope of the reserve curve is the price), so
the rebalancing terms cancel:
\[
  \frac{dV}{dy}
    \;=\; \underbrace{\frac{db_q}{dD}\frac{dD}{dy} + y\,\frac{dD}{dy}}_{=\,0}
       \;+\; D(y)
    \;=\; D(y) \;\geq\; 0.
\]
This is the envelope theorem: the first-order effect of the rebalancing
variable vanishes at the fair point, leaving the direct term. The
cancellation holds on the open band $(\alpha, \beta)$. Below it the pool
is all senior and $V(y) = y\,D(\alpha)$ rises linearly. Above it the
pool is all quote and $V$ is constant at $V_{\max}$. Both pieces are
non-decreasing and join the band continuously, so $V$ is non-decreasing
throughout.

The clamp now does its one job. With $\Dc(y) = \min(D(y), N)$ the
capped value $\widehat{V}$ has $\widehat{V}'(y) = \Dc(y) \leq N$
everywhere, so for $y > y'$,
\[
  0 \;\leq\; \widehat{V}(y) - \widehat{V}(y') \;\leq\; N\,(y - y').
\]
The candidate rate is $y = x/N$, so a unit step in $x$ moves
$f\widehat{V}$ by at most $f N \cdot \tfrac1N = f \leq 1$. A real
increase of at most $1$ changes a floor by $0$ or $1$, so
$P_A = \lfloor f \widehat{V}\rfloor$ rises by $0$ or $1$ per step, and
never falls.

\paragraph{Stage 2: $F$ is 1-Lipschitz.}
This is \cref{lem:F} below: $0 \leq F(j) - F(j') \leq j - j'$ for
$j > j'$.

\paragraph{Composition.}
The chain $x \mapsto x/N \mapsto P_A \mapsto F$ multiplies the per-stage
factors,
\[
  \underbrace{\frac{1}{N}}_{\div\, N}
  \;\cdot\;
  \underbrace{f N}_{P_A}
  \;\cdot\;
  \underbrace{1}_{F}
  \;=\; f \;\leq\; 1,
\]
since the junior owns at most the whole pool ($f \leq 1$). This product
is the \emph{loop gain} $g = f$: the dollars that return to the senior
NAV after one trip around the loop, per dollar sent in. Capping demand
at the supply holds it at or below one \emph{by construction}, with no
runtime condition to check.

On the integer lattice this gives the result directly. A unit step in
$x$ moves $P_A$ by $0$ or $1$ (Stage 1), and $F$ is 1-Lipschitz (Stage
2), so $F(P_A(x/N))$ rises by at most $1$. Subtracting the $-x$ term,
\[
  E(x{+}1) - E(x) \;\in\; \{-1,\ 0\}.
\]
$E$ is non-increasing, so the integers with $E \geq 0$ form an initial
segment $[0, x^\star]$. At its right endpoint the next step must be a drop of
$1$ (a flat step would keep $E \geq 0$), so $x^\star$ is the unique
integer with $E(x^\star) \geq 0 > E(x^\star{+}1)$. The drop being
exactly $1$ and $E$ integer-valued force $E(x^\star) = 0$, an exact
fixed point. Bisection returns this $x^\star$, the largest
self-consistent senior NAV. This proves Monotonicity.

\begin{lemma}[$F$ is non-decreasing and 1-Lipschitz in $\JTraw$]
\label{lem:F}
Fix $\STraw$ and the checkpoint and vary $j = \JTraw$. Then
$0 \leq F(j) - F(j') \leq j - j'$ for $j > j'$.
\end{lemma}

\begin{proof}
$F$ reaches $\STeff$ through two stages
(\cref{sec:overview:waterfall}), each with non-negative coefficients
$\in [0,1]$: attribution splits the raw delta by a claim fraction
$a \in [0,1]$ fixed at the checkpoint (each tranche's claim on a leg is
at most that leg's raw NAV, by conservation and non-negative effective
NAVs, so $a \leq 1$), and
settlement routes the result
through chains of $\min$, $\max$, addition and subtraction (slopes
$\{0,1\}$) and a yield split (fractions $w$, $1-w$ with $w$ capped).
So $\STeff$ is non-decreasing in $j$, and so is $\JTeff$. NAV
conservation finishes it: $\STeff + \JTeff = \STraw + j$, so the two
slopes sum to exactly $1$, and since both are non-negative each lies in
$[0,1]$. Thus $F$ is non-decreasing and 1-Lipschitz.
\end{proof}

\begin{remark}[The zero plateau and the loop gain]
\label{rem:plateau}
The clamp fixes the loop gain at $g = f \leq 1$
(\cref{sec:existence:mono}), which makes $E$ non-increasing but permits
it to be flat. Flatness can occur at zero whenever a floor step of $P_A$
passes through an $F$-branch of slope one, so the zero set
$\{x : E(x) = 0\}$ is a contiguous integer block at any gain, a
\emph{plateau} of self-consistent NAVs whenever it holds more than one
point. The gain bounds its length. Each step inside the block holds $E$
at $0$, so $F(P_A(x/N))$ rises by exactly $1$, and since $F$ is
1-Lipschitz and $P_A$ steps by $0$ or $1$ (\cref{sec:existence:mono}),
$P_A$ rises by exactly $1$ with it. A block of $k+1$ zeros thus forces
$f\widehat{V}$ across $k$ integer boundaries in $k$ steps of at most $f$
each. The start lies strictly less than one unit above its floor, so
crossing $k$ boundaries requires a total rise strictly above $k - 1$,
while the rise is at most $kf$. Hence $k(1-f) < 1$, and for $f < 1$ the
block holds at most $\lceil 1/(1-f) \rceil$ integers. At $f = 1$ the bound lapses
and the plateau is bounded only by the bracket itself. In real
arithmetic, with the floor in $P_A$ removed and $F$ read along its
affine branches (\cref{sec:overview:waterfall}), $E$ is continuous and
falls by at least $1-f$ per unit of $x$, so for every $f < 1$ it has a
unique real root. Removing the floor raises $E$ by less than one, so
that root lies in $[x^\star,\ x^\star + 1)$ and the published $x^\star$
is the nearest lattice point from below. What $f < 1$ buys is a plateau
bounded independently of the bracket and a real-arithmetic anchor, not a
singleton zero set. The clamp alone supplies neither.

The plateau costs nothing, because the published value is \emph{defined} to be
the largest plateau integer, the $x^\star$ of \cref{thm:main}, and that
choice is immaterial to the junior tranche. Across the block $E = 0$
gives $\STeff = x$, so by NAV conservation
$\JTeff = \STraw + P_A(x/N) - x$. A unit step inside the block forces
$P_A$ up by exactly $1$ (the step inference above), whence
\[
  \JTeff(x{+}1) - \JTeff(x)
    \;=\; \bigl(P_A((x{+}1)/N) - P_A(x/N)\bigr) - 1
    \;=\; 1 - 1 \;=\; 0.
\]
The junior's entitlement is identical at every plateau point. This
invariance is exactly as strong as $F$'s integer conservation: it relies
on $\STeff + \JTeff = \STraw + \JTraw$ holding at wei precision, the two
outputs not independently rounded, which the accountant guarantees by
construction (\cref{sec:overview:waterfall}). The senior's entitlement
is not invariant: $\STeff = x$ rises one wei per plateau point, so
self-consistent publications can differ by up to
$\lceil 1/(1-f) \rceil - 1$ wei when $f < 1$, and at $f = 1$ only the
bracket bounds the spread. Publishing the largest is one more tie
broken for the protected tranche (\cref{sec:setup:PA}). Existence does
not degrade at $f = 1$: the $+1$ slack in $x_{\max}$
(\cref{sec:existence:sign}) makes $E(x_{\max}) \leq -1 < 0$ whatever the
gain, so the bracket still straddles a crossing. A bracket-independent bound on the plateau needs
$f < 1$. Safety needs only the clamp.
\end{remark}

\subsection{Bracketing: the endpoints straddle the sign change}
\label{sec:existence:sign}

We exhibit a range whose endpoints have opposite signs by construction,
so no runtime sign-change check can fail.

\emph{Lower endpoint.} At $x = 0$, $E(0) = F(P_A(0)) = \STeff \geq 0$,
since the accountant holds effective NAVs non-negative, whatever the
junior collateral is worth. This consumes $\STraw \geq 0$: at $x = 0$
the conserved total is $\STraw + P_A(0) = \STraw$, so non-negative
outputs are consistent with conservation only when it holds (it also
keeps the bracket well-formed, $x_{\max} \geq 1$).

\emph{Upper endpoint.} $P_A$ is bounded above by $\lfloor f V_{\max}
\rfloor$. Since $V$ is
non-decreasing, $\widehat V \leq V \leq V_{\max}$ for every candidate
rate, whichever way the pool sorts its tokens. By NAV conservation
$\STeff = \STraw + \JTraw - \JTeff \leq \STraw
+ \JTraw \leq \STraw + \lfloor f V_{\max}\rfloor$ (using $\JTeff \geq
0$) for every candidate. At
$x_{\max} = \STraw + \lfloor f V_{\max}\rfloor + 1$ we therefore have
$E(x_{\max}) = \STeff - x_{\max} \leq -1 < 0$.

Both endpoints are computed once, at solve start, from frozen values.
The sign change is a theorem, not a checked precondition.

\subsection{What the argument assumes about the pool}
\label{sec:existence:generality}

The proof rests on structural facts an auditor verifies once per
market:

\begin{description}
  \item[Waterfall.] $F$ is non-decreasing and 1-Lipschitz in $\JTraw$
    (\cref{lem:F}), conserves NAV at wei precision with the two outputs
    not independently rounded, and returns non-negative effective NAVs.
    Bracketing consumes the last two directly
    (\cref{sec:existence:sign}), and \cref{lem:F}'s proof consumes them
    together with the yield split fraction $\in [0, 1]$. These are
    properties of the shared accountant and its yield model, and hold
    for any Royco market using the standard waterfall.
  \item[Pool.] The fair value $V$ is non-decreasing in $y$, with
    $V' = D \geq 0$ by the envelope theorem (\cref{sec:existence:mono}),
    the demand $D$ is non-increasing, and the all-quote corner value
    $V_{\max}$ is finite. Monotone $V$ holds for any constant-function
    pool. The finite corner is a property of the E-CLP's bounded band,
    and it is what makes the bracket endpoint $x_{\max}$ well-defined.
\end{description}

The proof never differentiates $F$, so it survives the
waterfall's kinks, and never inverts a pool invariant, so it transfers
unchanged across pool types: only the demand $D$ that the clamp caps is
pool-specific.

\paragraph{Out of scope.}
The solve fixes a manipulation-bounded rate. It does not monitor the
live pool's \emph{composition}. As senior shares accumulate and quote
drains, the published mark stays correct, since it reads the frozen
curve rather than live balances, but tracking that drift is a separate
operational concern. So is Royco's coverage and correlation parameter, a distinct
quantity from the pool's upper price bound $\beta$
(\cref{tbl:notation-math}). Neither is addressed by the fixed point.
